% !TEX spellcheck = en_US
% !TEX spellcheck = LaTeX
\documentclass[letterpaper,10pt,english]{article}
\usepackage{%
amsfonts,%
amsmath,%
amssymb,%
amsthm,%
babel,%
bbm,%
%biblatex,%
caption,%
centernot,%
color,%
enumerate,%
%enumitem,%
epsfig,%
epstopdf,%
etex,%
fancybox,%
framed,%
fullpage,%
%geometry,%
graphicx,%
hyperref,%
latexsym,%
mathptmx,%
mathtools,%
multicol,%
paralist,%
pgf,%
pgfplots,%
pgfplotstable,%
pgfpages,%
proof,%
psfrag,%
%subfigure,%
tcolorbox,%
tfrupee,%
tikz,%
times,%
%ulem,%
url,%
xcolor,%
mathpazo,%
}

\definecolor{shadecolor}{gray}{.95}%{rgb}{1,0,0}
\usepackage[margin=1in,top=0.75in]{geometry}
\usepackage[mathscr]{eucal}
\usepgflibrary{shapes}
\usepgfplotslibrary{fillbetween}
\usetikzlibrary{%
arrows,%
backgrounds,%
chains,%
decorations.pathmorphing,% /pgf/decoration/random steps | erste Graphik
decorations.text,% 
matrix,%
positioning,% wg. " of "
fit,%
patterns,%
petri,%
plotmarks,%
scopes,%
shadows,%
shapes.misc,% wg. rounded rectangle
shapes.arrows,%
shapes.callouts,%
shapes%
}

%\pgfplotsset{compat=newest} %<------ Here
\pgfplotsset{compat=1.11} %<------ Or use this one

\theoremstyle{plain}
\newtheorem{thm}{Theorem}[section]
\newtheorem{lem}[thm]{Lemma}
\newtheorem{prop}[thm]{Proposition}
\newtheorem{cor}[thm]{Corollary}
\newtheorem{clm}[thm]{Claim}

\theoremstyle{definition}
\newtheorem{defn}[thm]{Definition}
\newtheorem{conj}[thm]{Conjecture}
\newtheorem{exmp}[thm]{Example}
\newtheorem{axiom}[thm]{Axiom}
\newtheorem{assum}[thm]{Assumption}
\newtheorem{exerc}[thm]{Exercise}
\newtheorem*{defin}{Definition}

\theoremstyle{remark}
\newtheorem{rem}{Remark}
\newtheorem{note}{Note}

\newcommand{\iid}{\emph{i.i.d.}\ }
\newcommand{\Cov}{\operatorname{Cov}}
%\newcommand{\det}{\operatorname{det}}
%\newcommand{\Real}{\mathbb{R}}
\newcommand{\tr}{\operatorname{tr}}
\newcommand{\Var}{\operatorname{Var}}
\newcommand{\cov}{\operatorname{cov}}

%\renewcommand{\proof}[1]{\begin{proof}#1\end{proof}}
\newcommand{\EQ}[1]{\begin{equation*}#1\end{equation*}}
\newcommand{\EQN}[1]{\begin{equation}#1\end{equation}}
\newcommand{\eq}[1]{\begin{align*}#1\end{align*}}
\newcommand{\meq}[2]{\begin{xalignat*}{#1}#2\end{xalignat*}}
\newcommand{\norm}[1]{\left\lVert#1\right\rVert}
\newcommand{\inner}[1]{\left\langle #1\right\rangle}
\newcommand{\abs}[1]{\left\lvert#1\right\rvert}
\newcommand{\expect}[1]{\mathbb{E}\left[{#1}\right]}
\newcommand{\prob}[1]{\mathbb{P}\left[{#1}\right]}
\newcommand{\given}{\; \big\vert \;} 
\newcommand{\set}[1]{\left\{#1\right\}} 
\newcommand{\indicator}[1]{1_{\set{#1}}} 
\newcommand{\Ind}[1]{\mathbbm{1}_{#1}} 
\newcommand{\SetIn}[1]{\mathbbm{1}_{\set{#1}}} 
\newcommand{\red}[1]{\textcolor{red}{#1}} 
\newcommand{\floor}[1]{\left\lfloor#1\right\rfloor}
\newcommand{\ceil}[1]{\left\lceil#1\right\rceil}


\newcommand{\C}{\mathbb{C}}
\newcommand{\D}{\mathbb{D}}
\newcommand{\E}{\mathbb{E}}
\newcommand{\F}{\mathbb{F}}
\newcommand{\N}{\mathbb{N}}
\renewcommand{\P}{\mathbb{P}}
\newcommand{\Q}{\mathbb{Q}}
\newcommand{\R}{\mathbb{R}}
\newcommand{\Z}{\mathbb{Z}}

\newcommand{\bA}{\mathbf{A}}
\newcommand{\bI}{\mathbf{I}}
\newcommand{\bU}{\mathbf{1}}
\newcommand{\bx}{\mathbf{x}}
\newcommand{\bX}{\mathbf{X}}
\newcommand{\bZ}{\mathbf{Z}}


\newcommand{\cA}{\mathcal{A}}
\newcommand{\cB}{\mathcal{B}}
\newcommand{\cC}{\mathcal{C}}
\newcommand{\cF}{\mathcal{F}}
\newcommand{\cG}{\mathcal{G}}
\newcommand{\cM}{\mathcal{M}}
\newcommand{\cN}{\mathcal{N}}
\newcommand{\cP}{\mathcal{P}}
\newcommand{\cT}{\mathcal{T}}
\newcommand{\cX}{\mathcal{X}}
\newcommand{\cY}{\mathcal{Y}}

\newcommand{\sA}{\mathscr{A}}
\newcommand{\sB}{\mathscr{B}}
\newcommand{\sC}{\mathscr{C}}
\newcommand{\sE}{\mathscr{E}}
\newcommand{\sF}{\mathscr{F}}
\newcommand{\sG}{\mathscr{G}}
\newcommand{\sH}{\mathscr{H}}
\newcommand{\sL}{\mathscr{L}}
\newcommand{\sS}{\mathscr{S}}
\newcommand{\sT}{\mathscr{T}}
\newcommand{\sX}{\mathscr{X}}
\newcommand{\sY}{\mathscr{Y}}
\newcommand{\sZ}{\mathscr{Z}}

\newcommand{\tS}{\tilde{S}}
% Debug
\newcommand{\todo}[1]{\begin{color}{blue}{{\bf~[TODO:~#1]}}\end{color}}

% a few handy macros

\renewcommand{\le}{\leqslant}
\renewcommand{\ge}{\geqslant}
\newcommand\matlab{{\sc matlab}}
\newcommand{\goto}{\rightarrow}
\newcommand{\bigo}{{\mathcal O}}
%\newcommand{\half}{\frac{1}{2}}
%\newcommand\implies{\quad\Longrightarrow\quad}
\newcommand\reals{{{\rm l} \kern -.15em {\rm R} }}
\newcommand\complex{{\raisebox{.043ex}{\rule{0.07em}{1.56ex}} \hskip -.35em {\rm C}}}


% macros for matrices/vectors:

% matrix environment for vectors or matrices where elements are centered
\newenvironment{mat}{\left[\begin{array}{ccccccccccccccc}}{\end{array}\right]}
\newcommand\bcm{\begin{mat}}
\newcommand\ecm{\end{mat}}

% matrix environment for vectors or matrices where elements are right justifvied
\newenvironment{rmat}{\left[\begin{array}{rrrrrrrrrrrrr}}{\end{array}\right]}
\newcommand\brm{\begin{rmat}}
\newcommand\erm{\end{rmat}}

% for left brace and a set of choices
%\newenvironment{choices}{\left\{ \begin{array}{ll}}{\end{array}\right.}
\newcommand\when{&\text{if~}}
\newcommand\otherwise{&\text{otherwise}}
% sample usage:
%\delta_{ij} = \begin{choices} 1 \when i=j, \\ 0 \otherwise \end{choices}


% for labeling and referencing equations:
\newcommand{\eql}{\begin{equation}\label}
\newcommand{\eqn}[1]{(\ref{#1})}
% can then do
%\eql{eqnlabel}
%...
%\end{equation}
% and refer to it as equation \eqn{eqnlabel}.
% some useful macros for finite difference methods:
\newcommand\unp{U^{n+1}}
\newcommand\unm{U^{n-1}}
% for chemical reactions:
\newcommand{\react}[1]{\stackrel{K_{#1}}{\rightarrow}}
\newcommand{\reactb}[2]{\stackrel{K_{#1}}{~\stackrel{\rightleftharpoons}
 {\scriptstyle K_{#2}}}~}
\makeatletter
\def\th@plain{%
\thm@notefont{}% same as heading font
\itshape % body font
}
\def\th@definition{%
\thm@notefont{}% same as heading font
\normalfont % body font
}
\makeatother
\date{}


\title{Lecture-06: Transformation of random vectors}
\author{}

\begin{document}
\maketitle

\section{Functions of random variables}
\begin{defn}
Borel measurable sets on a space $\R^n$ is denoted by $\sB(\R^n)$ and generated by the collection $(\pi_i^{-1}(-\infty,x]: x \in \R, i \in [n])$. 
A function $g: \R^n\to \R^m$ is called \textbf{Borel measurable} function, 
if $g^{-1}(B_m) \in \sB(\R^n)$ for any $B_m \in \sB(\R^m)$. 
\end{defn}
\begin{prop}
Consider a random variable $X: \Omega \to \R$ defined on the probability space $(\Omega, \sF, P)$. 
Suppose $g: \R \to \R$ is function such that $g^{-1}(-\infty, x] \in \sB(\R)$, 
then $g(X)$ is a random variable. 
\end{prop}
\begin{proof}
We represent $g(X)$ by a map $Y: \Omega \to \R$ such that $Y(\omega) \triangleq (g \circ X)(\omega)$ for all outcomes $\omega \in \Omega$. 
We further check that for any half open set $B_x = (-\infty,x]$, we have $Y^{-1}(B_x) = (X^{-1}\circ g^{-1})(B_x)$. 
Since $g^{-1}(B_x) \in \sB(\R)$, it follows that $Y^{-1}(B_x) \in \sF$ by the definition of random variables.  
\end{proof}
\begin{shaded}
\begin{exmp}[Monotone function of random variables] 
Let $g:\R \to \R$ be a monotonically increasing function, 
then $g^{-1}(-\infty, x] %= (g^{-1}(-\infty), g^{-1}(x)] 
\in \sB(\R)$ for all $x \in \R$. 
Consider a random variable $X: \Omega \to \R$ defined on the probability space $(\Omega, \sF, P)$, 
then $Y \triangleq g(X)$ is a random variable with distribution function
\EQ{
F_Y(y) = P\set{g(X) \le y} = P\set{X \le g^{-1}(y)} = F_{X}(g^{-1}(y)). 
}
Here, $g^{-1}(y)$ is the functional inverse, and not inverse image as we have been seeing typically. 
We can think $g^{-1}(y) = g^{-1}\set{y}$, though this inverse image has at most a single element since $g$ is monotonically increasing. 
\end{exmp}
\end{shaded}
\begin{shaded}
\begin{exmp}
Consider a positive random variable $X: \Omega \to \R_+$ defined on a probability space $(\Omega,\sF,P)$.    
Let $g: \R_+ \to \R_+$ be such that $g(x) = e^{-\theta x}$ for all $x \in \R_+$ and some $\theta > 0$.  
Then, $g$ is monotonically decreasing in $X$ and $x = g^{-1}(y) = -\frac{1}{\theta}\ln y$. 
This implies that $g^{-1}(-\infty, y] = [-\frac{1}{\theta}\ln y, \infty) \in \sB(\R_+)$ for all $y \in \R_+$. 
Thus $g$ is a measurable function, and $Y = g(X)$ is a random variable.  
\end{exmp}
\end{shaded}
%\begin{shaded}
\begin{prop}[Independence of function of random variables] 
Let $g:\R \to \R$ and $h: \R \to \R$ be functions such that $g^{-1}(-\infty, x]$ and $h^{-1}(-\infty, x]$ are Borel sets for all $x \in \R$. 
Consider independent random variables $X$ and $Y$ defined on the probability space $(\Omega, \sF, P)$, 
then $g(X)$ and $h(Y)$ are independent random variables. 
\end{prop}
\begin{proof}
For any $u,v \in \R$, we can define inverse images $A_g(u) \triangleq g^{-1}(-\infty,u]$ and $A_h(v) \triangleq h^{-1}(\infty, v]$. 
Since $g,h$ are Borel measurable, we have $A_g(u), A_h(v) \in \sB(\R)$.   
We can write the following outcome set equality for the joint event 
\EQ{
\set{g(X) \le u}\cap\set{h(Y) \le v} = \set{X \in g^{-1}(-\infty, u]}\cap\set{Y \in h^{-1}(-\infty,v]} = X^{-1}(A_g(u))\cap Y^{-1}(A_h(v)) \in \sF.
}
Since $X$ and $Y$ are independent random variables, it follows that $X^{-1}(A_g(u))$ and $Y^{-1}(A_h(v))$ are independent events, 
and the result follows. 
\end{proof}
%\end{shaded}

\section{Function of random vectors}
\begin{prop}
Consider a random vector $X: \Omega \to \R^n$ defined on the probability space $(\Omega, \sF, P)$, 
and a Borel measurable function $g: \R^n \to \R^m$ such that $A_g(y) \triangleq \cap_{j=1}^m\set{x \in \R^n: g_j(x) \le y_j} \in \sB(\R^n)$ for all $y \in \R^m$. 
Then, $g(X): \Omega \to \R^m$ is a random vector. 
The joint distribution function $F_Y: \R^m \to [0,1]$ for the vector $Y\triangleq g(X)$ is given by 
\EQ{
F_Y(y) = P(X^{-1}(A_g(y))),\quad \text{ for all } y \in \R^m.
} 
\end{prop}
\begin{shaded}
\begin{exmp}[Sum of random variables]
For a random vector $X: \Omega \to \R^n$ defined on a probability space $(\Omega, \sF, P)$. 
Define an addition function $+: \R^n \to \R$ such that $+(x) = \sum_{i=1}^nx_i$ for any $x \in \R^n$. 
We can verify that $+$ is a Borel measurable function and hence $Y = +(X) = \sum_{i=1}^nX_i$ is a random variable. 
When $n=2$ and $X$ is a continuous random vector with density $f_X: \R^2 \to \R_+$, we can write 
\EQ{
F_Y(y) 
= P(\set{Y \le y}) 
= P(\set{X_1+X_2 \le y}) 
%= \int_{x_1 + x_2 \le y}f_X(x_1, x_2)dx_1dx_2
= \int_{x_1 \in \R}\int_{x_2 \le y-x_1}f_X(x_1,x_2)dx_1dx_2.
}
By applying a change of variable $(x_1, t)= (x_1, x_1 + x_2)$ and changing the order of integration, 
we see that 
\EQ{
F_Y(y) = \int_{t \le y}dt \int_{x_1 \in \R} dx_1 f_X(x_1,t-x_1). 
} 
When $Y$ is a continuous random vector, we can write 
\EQ{
f_Y(y) = \frac{dF_Y(y)}{dy} = \int_{x\in \R}f_X(x, y-x)dx.
}
When $X: \Omega \to \R^2$ is an independent vector, then $f_X(x) = f_{X_1}(x_1)f_X(x_2)$ for all $x\in\R^2$. 
Therefore, the density of the sum $X_1+X_2$ is given by  
\EQ{
f_Y(y) = \int_{x \in \R}dx f_{X_1}(x)f_{X_2}(y-x) = (f_{X_1}\ast f_{X_2})(y), 
}
where $\ast: \R^{\R} \times \R^{\R} \to \R^{\R}$ is the convolution operator. 
\end{exmp}
\end{shaded}
\begin{thm}
For a continuous random vector $X:\Omega \to \R^m$ defined on the probability space $(\Omega, \sF, P)$ with density $f_X: \R^m \to \R_+$ and an injective and smooth Borel measurable function $g: \R^m \to \R^m$, 
such that $Y=g(X)$ is a continuous random vector. 
Then the density of random vector $Y$ is given by 
\EQ{
f_Y(y) = \frac{f_X(x)}{\abs{J(y)}},
} 
where $x=g^{-1}(y)$ and $J(y)= (J_{ij}(y) \triangleq \frac{\partial y_j}{\partial x_i}: i,j\in [m])$ is the Jacobian matrix.  
\end{thm} 
\begin{proof}
For an injective map $g: \R^m \to \R^m$ we have $\set{x} = g^{-1}\set{y}$ for any $y \in g(\R^m)$. 
Further, since $g$ is smooth, we have $dy = J(y)dx + o(\abs{dx})$, 
and thus
\EQN{
\label{eqn:Jacobian}
\abs{dy} = \abs{J(y)}\abs{dx} + o(\abs{dx}).
} 
Defining set $dB(y) \triangleq \set{w \in \R^m: y_j \le w_j \le y_j+dy_j}$, 
we observe that for any continuous random vector $Y: \Omega \to \R^m$, we have 
\EQN{
\label{eqn:density}
P \circ Y^{-1} \left(dB(y)\right)= f_Y(y)\abs{dy} = P\circ X^{-1}\Big(dB(x)\Big) 
= f_X(x)\abs{dx}. 
}
We get the result by combining~\eqref{eqn:Jacobian} and~\eqref{eqn:density}.
\end{proof}
\begin{shaded}
\begin{exmp}[Sum of random variables]
Suppose that $X: \Omega \to \R^2$ is a continuous random vector and $Y_1 = X_1 + X_2$. 
Let us compute $f_{Y_1}(y_1)$ using the above theorem. 
Let us define a random vector $Y: \Omega \to \R^2$ such that $Y = (X_1+X_2, X_2)$ 
so that $\abs{J(y)} = 1$.
This implies, $f_{Y}(y) = f_{X}(x)$.
Thus, we may compute the marginal density of $Y_1$ as,
\EQ{
f_{Y_1}(y_1) 
= \int_{-\infty}^{\infty} f_X(x)\SetIn{x_2= y_2, x_1+x_2 = y_1}dy_2 
= \int_{-\infty}^{\infty} f_X(y_1 -y_2,y_2)dy_2.
}
If $X$ is an independent random vector, then 
%\textbf{Special Case} : Suppose $X_1$ and $X_2$ are independent.
%Then, 
\EQ{
f_{Y_1}(y_1) = \int_{-\infty}^{\infty} f_{X_1}(y_1-y_2)f_{X_2}(y_2)dy_2 = (f_{X_1} * f_{X_2})(y_1),
}
where $*$ represents convolution.
\end{exmp}
\end{shaded}
%Now, let us try to compute the above density $f_{Y_1}{y_1}$ in a more direct manner.
%We know,
%\EQ{
%f_{Y_1}(y_1) = \frac{dF_{Y_1}(y_1)}{dy_1}.
%}
%where, $F_{Y_1}(y_1) = P\set{Y_1 \le y_1} = P\set{X_1 + X_2 \le y_1}$. 
%Representing this probability as the area under the joint density function, we have,
%\EQ{
%F_{Y_1}(y_1) = \int_{\infty}^{\infty}dx_1 \int_{-\infty}^{y_1 - x_1} f_{X_1,X_2}(x_1,x_2)dx_2.
%}
%Now, by applying change of variable $(x_1, t)= (x_1, x_1 + x_2)$ and changing the order of integration, we see that,
%\EQ{
%F_{y_1}(y_1) = \int_{-\infty}^{y_1}dt \int_{-\infty}^{\infty} dx_1 f_{X_1,X_2}(x_1,t-x_1).
%} 
%Finally, we get the expected result
%\EQ{
%f_{Y_1}(y_1) = \frac{dF_{Y_1}(y_1)}{dy_1} = \int_{-\infty}^{\infty}  f_{X_1,X_2}(x_1,y_1-x_1)dx_1.
%}
%\section{Characteristic function}
%Suppose that $X$ is a continuous random variable with density $f_X(x)$. Then, recall that the characteristic function is given by 
%\EQ{
%\E[e^{juX}] = \int_{-\infty}^{\infty}e^{juX}f_X(x)dx, 
%}
%where, $e^{juX} = \cos{uX} + j\sin{uX}$.
\end{document}
